\section{R1. Scalar companion and risk
proof}\label{r1.-scalar-companion-and-risk-proof}\label{sup:R1}

This is a new statement for the revision, with the complete argument
below. It does not alter the confidence interval or assert pointwise
minimax squared risk.

Let \(Y=\psi+\tau g\), \(g\sim N(0,1)\), and suppose the supplied
structural range is \(|\psi|\le h\). Let \(I\) be the intersection of
any candidate interval with \([Y-d,Y+d]\) and \([-h,h]\), where
\(d=z_{1-\alpha/4}\tau\). The candidate may depend arbitrarily on the
data. The reported confidence interval may retain its existing
empty-intersection convention \(\{0\}\).

Define a \textbf{separate scalar estimator} by \(\widehat\psi=0\) when
\(h\le\tau\). Otherwise, use the midpoint of the actual nonempty
intersection; on an empty intersection use \(\Pi_{[-h,h]}Y\). Replacing
an exactly zero singleton by the clipped direct estimate also satisfies
the bound. Then, for every admissible mean,

\[E(\widehat\psi-\psi)^2\le
\begin{cases}
h^2,&h\le\tau,\\
\tau^2\{1+2z_{1-\alpha/4}\sqrt{2/\pi}+z_{1-\alpha/4}^2\},&h>\tau.
\end{cases}\]

Thus the risk is at most \(C_\alpha\min(h^2,\tau^2)\), uniformly over
the norm budget and signal magnitude. No candidate coverage, separation,
rank recovery, or independence from \(Y\) is needed for this
\textbf{risk} bound. Honesty of the reported interval continues to
require its original assumptions and calibration.

\textbf{Proof.} On \(h\le\tau\), zero has squared error at most \(h^2\).
Otherwise, on a nonempty intersection its midpoint \(m\) is in
\([Y-d,Y+d]\), so \(|m-\psi|\le|Y-\psi|+d\). On an empty intersection,
projection onto an interval containing \(\psi\) is nonexpansive relative
to \(\psi\): \(|\Pi_{[-h,h]}Y-\psi|\le|Y-\psi|\). The same argument
covers a zero-singleton replacement. Consequently the error on every
observation is at most \(\tau|g|+d\). Taking expectations, using
\(Eg^2=1\) and \(E|g|=\sqrt{2/\pi}\), proves the second branch. Since
\(C_\alpha\ge1\), both branches imply the final inequality. No
conditioning on a successful pilot or deletion of rare events occurs.

For A4/P1 and the normalized diagonal loading, \(h=M\sqrt{r/p}\). The
cap-valid intervals already intersect the direct interval at error
\(\alpha/2\), so \(d=z_{1-\alpha/4}\tau\) is exactly their existing
direct radius. The scalar companion therefore removes the \(M^2\) loss
from the arbitrary zero sentinel without changing any coverage event or
interval length. For general fixed \(L\), replace \(\tau\) by
\(\tau\|L\|_F\) and use any valid supplied structural range. This is a
noise-scale risk bound, not a proof of the smaller P1 pointwise
squared-risk rate or of the ideal full-data minimax risk.
